\answer{ex:22:1}{(а)~$32\%$ и $100\%$. (б)~$R_\uparrow-R_\downarrow\ge2\sqrt{2000\,\text{Гц}/1\,\text{с}}\approx89$~Гц, около $4.5\%$ суммарной частоты.}
\answer{ex:22:2}{(а)~$n\ge4(p_1q_1+p_2q_2)/\Delta^2\approx560$. (б)~$\approx56$. (в)~Удары одной культуры коррелированы через её настройку; единица статистики --- культура.}
\answer{ex:22:3}{(а)~Детальный баланс $\rho P K$ симметричен; единственно при неприводимом $K$. (б)~Доля промахов $=1/\langle1/P\rangle\le\langle P\rangle$, равенство при $P=\text{const}$. (в)~$0.18$ против $0.5$. (г)~Такие настройки поглощающие; вся масса собирается в них.}
\answer{ex:22:4}{(б)~Параллелограмм площади $4h^2=0.04$ (с учётом ограничения $p$ --- численно $0.045$). (в)~$\approx0.18$, в согласии с $0.19$ на рис.~\ref{fig:dishmodel}.}
\answer{ex:22:5}{(а)~$0.49$; в поглощающей области $44\%$ культур (после $2000$ розыгрышей --- $32\%$); остальные уходят в область, где промах почти всегда. (б)~$0.80$ и $1.00$. (в)~$0.25$, $0.32$, $0.38$, $0.36$, $0.32$, $0.23$ при $\sigma=0.02$--$0.64$ ($\pm0.03$); оптимум при $\sigma\approx0.1\approx h$.}
\answer{ex:22:6}{(а)~$\ge5$ уровней; $8$ электродов хватает. (б)~Нет: нужно $r_{\max}\ge25/T=100$~Гц. (в)~$\sqrt{10}\approx3$ уровня, $\approx1.7$~бита; высоту --- местом, расстояние --- частотой.}
\answer{ex:22:7}{Открытая задача. Время случайного поиска растёт с $d$ примерно как отношение объёмов $(L/h)^d$, у поиска со знаком ошибки --- линейно по $d$; опыт с перестановкой предсказуемости и силы стимулов разделяет гипотезы, если в каждой группе несколько десятков культур.}
